\begin{figure}[!hbtp]
    \centering
    \resizebox{0.75\columnwidth}{!}{%
    \begin{tikzpicture}[
        recntagle/.style={
            rectangle, rounded corners=3pt, inner sep=5pt, align=center,
            font=\small\bfseries, draw=black!70, fill=blue!10, line width=0.8pt
        },
        title_block/.style={
            rectangle, rounded corners=3pt, inner sep=4pt, align=center,
            font=\small\bfseries, draw=black!70, fill=blue!10, line width=0.8pt
        },
        arrow/.style = {thick, ->, >={Stealth[length=2.5mm]}},
        mini_arrow/.style = {semithick, ->, >={Stealth[length=1.8mm]}},
        box_delimitador/.style = {
            rectangle, rounded corners=8pt, draw=black!70, fill=black!2, line width=0.8pt
        },
        sub_block/.style = {
            rectangle, rounded corners=3pt, inner sep=4pt, align=center,
            draw=black!70, fill=blue!5, line width=0.7pt, font=\small
        },
        sub_item/.style = {
            rectangle, rounded corners=2.5pt, inner sep=3pt, align=center,
            draw=black!70, fill=white, line width=0.6pt, font=\scriptsize
        },
        emb_item/.style = {
            rectangle, rounded corners=2.5pt, inner sep=3pt, align=center,
            draw=black!70, line width=0.6pt, font=\scriptsize
        },
        emb/.style = {
            rectangle, rounded corners=2pt, draw=black!70, line width=0.6pt, minimum width=1.45cm,
            minimum height=0.7cm, font=\scriptsize, inner sep=2pt, align=center
        }
    ]

        % 1. Datasets
        \node (dataset) [recntagle, minimum width=6.2cm] at (0,0) {1. \textit{Datasets} histopatológicos};

        % 2. SSN: Modelagem em grafo e extração de características (via topológica)
        \node (ssn_box) [box_delimitador, minimum width=5.8cm, minimum height=2.5cm] at (-3.6,-3.5) {};
        \node (ssn_title) [title_block] at (ssn_box.north) {2. Modelagem Topológica (SSN)};
        
        \node (ssn_img) [sub_item, minimum width=1.1cm, minimum height=1.1cm] at (-5.2, -3.55) {Imagem\\original};
        
        \node (ssn_graph) [sub_item, minimum width=1.2cm, minimum height=1.1cm] at (-3.6, -3.55) {};
        \node [font=\scriptsize, align=center] at (-3.6, -3.22) {Grafo};
        \filldraw [blue!70!black] (-3.95, -3.75) circle (1.8pt);
        \filldraw [blue!70!black] (-3.95, -3.45) circle (1.8pt);
        \filldraw [blue!70!black] (-3.25, -3.75) circle (1.8pt);
        \filldraw [blue!70!black] (-3.25, -3.45) circle (1.8pt);
        \draw [black!60, line width=0.5pt] (-3.95, -3.75) -- (-3.95, -3.45) -- (-3.25, -3.75) -- (-3.25, -3.45) -- (-3.95, -3.75);
        \draw [black!60, line width=0.5pt] (-3.95, -3.45) -- (-3.25, -3.45);
        
        \node (ssn_vec) [sub_item, minimum width=1.1cm, minimum height=1.1cm] at (-2.0, -3.55) {Vetor\\$X_{\mathrm{top}}$};
        
        \draw [mini_arrow] (ssn_img.east) -- (ssn_graph.west);
        \draw [mini_arrow] (ssn_graph.east) -- (ssn_vec.west);

        % 3. RNN: Vetor de entrada -> RNN -> Embedding (via topológica)
        \node (rnn_box) [box_delimitador, minimum width=5.8cm, minimum height=2.5cm] at (-3.6,-7.0) {};
        \node (rnn_title) [title_block] at (rnn_box.north) {3. Codificador Recorrente (RNN)};
        
        \node (rnn_vec_in) [sub_item, minimum width=1.1cm, minimum height=1.1cm] at (-5.2, -7.05) {Vetor\\$X_{\mathrm{top}}$};
        
        \node (rnn_unit) [sub_item, minimum width=1.2cm, minimum height=1.1cm] at (-3.6, -7.05) {Célula\\RNN};
        \draw [->, >={Stealth[length=2mm]}, line width=0.7pt, black!75] ([xshift=4pt]rnn_unit.north) to [out=65, in=115, looseness=3.5] ([xshift=-4pt]rnn_unit.north);
        
        \node (rnn_emb_out) [emb_item, fill=blue!15, minimum width=1.1cm, minimum height=1.1cm] at (-2.0, -7.05) {\textit{Embedding}\\$z_{\mathrm{top}}$};
        
        \draw [mini_arrow] (rnn_vec_in.east) -- (rnn_unit.west);
        \draw [mini_arrow] (rnn_unit.east) -- (rnn_emb_out.west);

        % 4. Pré-processamento visual (via visual - bloco superior)
        \node (prep_box) [box_delimitador, minimum width=5.8cm, minimum height=2.5cm] at (3.6,-3.5) {};
        \node (prep_title) [title_block] at (prep_box.north) {4. Pré-processamento Visual};
        
        \node (prep_resize) [sub_item, minimum width=1.7cm, minimum height=1.1cm] at (2.2, -3.55) {Redimensionar\\$256 \times 256$};
        \node (prep_aug) [sub_item, text width=2.1cm, minimum height=1.1cm] at (5.0, -3.55) {\textit{Data Augmentation}\\(rotações)};
        \draw [mini_arrow] (prep_resize.east) -- (prep_aug.west);

        % 5. Modelos de Visão Computacional (via visual - bloco inferior)
        \node (vis_box) [box_delimitador, minimum width=5.8cm, minimum height=2.5cm] at (3.6,-7.0) {};
        \node (vis_title) [title_block] at (vis_box.north) {5. Modelos Visuais};
        
        \node (vis_models) [sub_item, minimum width=1.8cm, minimum height=1.1cm] at (2.25, -7.05) {Backbone\\visual};
        \node (vis_emb_out) [emb_item, fill=orange!20, minimum width=1.5cm, minimum height=1.1cm] at (4.95, -7.05) {\textit{Embedding}\\$\tilde{z}_{\mathrm{vis}}$};
        \draw [mini_arrow] (vis_models.east) -- (vis_emb_out.west);

        % 6. MLP: Concatenação + Classificação
        \node (mlp_box) [box_delimitador, minimum width=9.8cm, minimum height=3.0cm] at (0,-10.9) {};
        \node (mlp_title) [title_block] at (mlp_box.north) {6. Fusão e Classificação (MLP)};
        
        \node [font=\footnotesize, align=center] at (-2.6,-10.15) {Concatenação de\\\textit{embeddings}};
        \node (cat_rnn) [emb, fill=blue!15]   at (-3.4,-11.25) {Topológico\\$z_{\mathrm{top}}$};
        \node (cat_vis) [emb, fill=orange!20] at (-1.8,-11.25) {Visual\\$\tilde{z}_{\mathrm{vis}}$};
        
        \node (mlp_sub) [sub_block, minimum width=2.8cm, minimum height=1.1cm] at (2.5,-11.25) {Camadas densas\\(MLP + \textit{Dropout})};
        \draw [arrow] ([xshift=0.15cm]cat_vis.east) -- (mlp_sub.west)
            node [midway, above, font=\scriptsize] {$z_{\mathrm{fus}}$};

        % 7. Métricas
        \node (metrics) [recntagle, minimum width=6.2cm] at (0,-13.8) {7. Extração de métricas};

        % Conexões principais entre os blocos
        \draw [arrow] (dataset.south) -- ++(0,-0.6) -| (ssn_title.north);
        \draw [arrow] (dataset.south) -- ++(0,-0.6) -| (prep_title.north);
        
        \draw [arrow] (ssn_box.south) -- (rnn_title.north)
            node [midway, right, font=\footnotesize, xshift=1pt] {$X_{\mathrm{top}} \in \mathbb{R}^{15\times 18}$};
            
        \draw [arrow] (prep_box.south) -- (vis_title.north)
            node [midway, right, font=\footnotesize, xshift=1pt] {Imagem $256\times256$};
            
        \draw [arrow] (rnn_box.south) -- (rnn_box.south |- mlp_box.north)
            node [midway, left, font=\footnotesize, xshift=-1pt] {$z_{\mathrm{top}}$};
            
        \draw [arrow] (vis_box.south) -- (vis_box.south |- mlp_box.north)
            node [midway, right, font=\footnotesize, xshift=1pt] {$\tilde{z}_{\mathrm{vis}}$};
            
        \draw [arrow] (mlp_box.south) -- (metrics.north)
            node [midway, right, font=\footnotesize, align=left, xshift=2pt] {Classificação\\histopatológica};

    \end{tikzpicture}%
    }
    \caption{Pipeline geral da abordagem proposta: os \textit{datasets} histopatológicos alimentam duas vias paralelas. Na via topológica, a imagem original é modelada como um grafo complexo (SSN), gerando a sequência de características $X_{\mathrm{top}}$, que é processada por um codificador recorrente (RNN) para originar o \textit{embedding} topológico $z_{\mathrm{top}}$. Na via visual, a imagem passa por pré-processamento (redimensionamento para $256\times256$ e aumento de dados por rotação) antes de ser processada pelo \textit{backbone} visual para gerar o \textit{embedding} visual $\tilde{z}_{\mathrm{vis}}$. Ambos os \textit{embeddings} são equalizados, concatenados e classificados por um MLP, gerando as predições finais para extração das métricas.}
    \label{fig:pipeline-geral}
    \label{fig:pipeline_tcc}
\end{figure}
